\documentclass[11pt]{article}
\usepackage[%
  papersize={160mm,120mm},
  bottom=2mm,
  top=2mm,
  textwidth=150mm,
  textheight=115mm,
  hcentering,
]{geometry}

% Mathtools with AMS Math
\usepackage{mathtools}

% Use TeX Gyre Schola font family with Fourier math fonts
\usepackage[T1]{fontenc}
\usepackage{tgschola}
\usepackage{fouriernc}
\usepackage{parskip}
\setlength{\parskip}{0.25em}
\renewcommand{\baselinestretch}{0.9}

\usepackage{enumitem}
\setlist{itemsep=.05em,parsep=.25em}

\usepackage{graphicx}
\usepackage{amsmath, caption, fancybox, soul,xspace}
\usepackage[dvipsnames]{xcolor}
\usepackage{hyperref}
\hypersetup{colorlinks,linkcolor=,urlcolor=blue}

\usepackage{axodraw2}

\usepackage{fancyvrb}

\usepackage{minibox}

\newcommand{\be}{\begin{equation}}
\newcommand{\bi}{\begin{itemize}}
\newcommand{\bn}{\begin{enumerate}}
\newcommand{\bv}{\begin{Verbatim}}
\newcommand{\ee}{\end{equation}}
\newcommand{\ei}{\end{itemize}}
\newcommand{\en}{\end{enumerate}}
\newcommand{\ev}{\end{Verbatim}}

\newcommand{\lb}{\left(}
\newcommand{\rb}{\right)}

\newcommand{\fire}{\texttt{FIRE}\xspace}
\newcommand{\fuchsia}{\texttt{Fuchsia}\xspace}
\newcommand{\hyperint}{\texttt{HyperInt}\xspace}
\newcommand{\litered}{\texttt{LiteRed}\xspace}
\newcommand{\reduze}{\texttt{Reduze2}\xspace}

\newcommand{\blue}[1]{{\color{blue} #1}}

\newcommand{\x}{\;}
\newcommand{\vs}{\vspace{2mm}}

\newcommand{\abs}[1]{\left|#1\right|}
\newcommand{\D}{\mathrm{d}}
\newcommand{\eps}{\epsilon}
\newcommand{\linux}{\texttt{Linux}\xspace}
\newcommand{\maple}{\texttt{Maple}\xspace}
\newcommand{\maxima}{\texttt{Maxima}\xspace}
\newcommand{\maximasage}{\texttt{Maxima/Sage}\xspace}
\newcommand{\python}{\texttt{Python}\xspace}
\newcommand{\rank}{\mathrm{rank}}
\newcommand{\sage}{\texttt{SageMath}\xspace}
\newcommand{\code}[1]{\texttt{#1}}
\newcommand{\F}[1]{\texttt{#1}} % use this to style function names
\newcommand{\V}[1]{\mathbf{#1}} % use this to style vector names
\newcommand{\prompt}[2]{\textcolor{prompt}{#1} \textcolor{command}{#2}}
\newcommand{\functionitem}[2]{\item[$\F{#1}(#2)$\hfill\textit{(function)}]}
\newcommand{\classitem}[1]{\item[$#1$\hfill\textit{(class)}]}


\newcommand{\refx}[1]{\href{1}{\scriptsize #1}}
\newcommand{\titlebox}[1]{\centerline{\color{black}\shadowbox{\begin{Bcenter}{\Large \color{Blue} #1}\end{Bcenter}}}}
\newcommand{\titleboxx}[2]{\centerline{\color{black}\shadowbox{\begin{Bcenter}\Large\color{black} #1\vspace{1mm}\\  {\color{white}q}{\color{black}#2}{\color{white}q}\end{Bcenter}}}}

\newcommand{\bbox}[1]{\color{black}\boxed{{\color{Plum}#1}}}

%\setmonofont{Go-Mono}
\definecolor{textcolor}{RGB}{60,60,60}

\newcommand{\spr}[2]{#1\!\cdot\!#2}

\everymath{\color{Plum}}
\everydisplay{\color{Plum}}

\setlength{\parindent}{0mm}

\begin{document}

\color{textcolor}

\begin{center}
  \mbox{}
  \vfill
  {\LARGE \bf \scshape Introduction \\ \Large to \LARGE \\ Differential Equations \\ \Large for \LARGE \\ Feynman Integrals \\}
  \vfill
  {\large Oleksandr Gituliar} \\ \href{mailto:oleksandr@gituliar.net}{oleksandr@gituliar.net}
  \vfill
  \includegraphics[height=1.2cm]{img/logo_uhh}\\
  II. Institut f\"ur Theoretische Physik \\ Universit\"at Hamburg
  \vfill
  {Computer Algebra in Particle Physics 2017 \\ DESY (Hamburg)}
  \vfill
  \mbox{}
\end{center}
\newpage
 

\titlebox{Plan for Today}
\begin{enumerate}
  \item Methods for caluclating Feynman integrals (20')
  \begin{itemize}
    \item Integration by Parts (IBP)
    \item Method of Differential Equations (MDE)
  \end{itemize}
  \item MDE \& Solution Techniques (30')
  \begin{itemize}
    \item ordinary vs partial
    \item epsilon form
    \item boundary conditions
  \end{itemize}
  \item Examples (30')
  \item MDE \& Limitations (10')
\end{enumerate}
\newpage

\titlebox{I.1 Feynman Integrals Calculus}
\bi \item Feynman/Schwinger parameters
          \bi \item \hyperint and hyperlogarithms
              \ei
    \item Mellin-Barnes (MB) representation
          \bi \item popular for massive diagrams
              \ei
    \item Integration-by-parts (IBP) reduction
          \bi \item universal and easy-to-use
              \item \fire, \litered, \reduze
              \ei
    \item Method of differential equations (MDE)
          \bi \item huge potential
              \ei
    \ei
\newpage

\titlebox{I.2 Method of Differential Equations}
\bn
  \item Find IBP rules
    \bi
      \item highly automated
      \item \fire, \litered, \reduze
    \ei
  \item Derive ODE
  \item Find $\eps$-form
    \bi
      \item already very powerful (less-understood than IBP)
      \item Lee method: \fuchsia (2016), \texttt{epsilon} (2017)
      \item other:
    \ei
  \item Solve ODE
  \item Find constants of integration
    \bi
      \item depends on the problem
      \item may be the most complex task
    \ei
\en
\newpage

\titleboxx{\Large Example I}{\Large \bf One-Loop Massive Self-Energy}
%{\small
%{\bf Scalar integrals} are the most {\em convenient} to compute.
%Therefore, we assume that integrals with vector or spinor indices (e.g. amplitudes) are already expressed as a combination of scalar integrals and constant objects which reveal the index structure.}
\bi
  \item {\bf Functions}: {\em from vectors to scalars}
$$ \raisebox{-28pt}{
     \SetScale{0.5}
     \begin{axopicture}(200,120)
       \SetOffset(0,10)
       \DoubleArc[arrow](100,50)(40,0,180){2}
       \Text(100,105){$l$}
       \DoubleArc[arrow](100,50)(40,180,360){2}
       \Text(100,-5){$l-p$}
       \Text(30,70){$p$}
       \Gluon(0,50)(60,50){5}{4}
       \Vertex(60,50){2}
       \Text(170,70){$p$}
       \Gluon(140,50)(200,50){5}{4}
       \Vertex(140,50){2}
     \end{axopicture}}
   =
 \Pi^{\mu\nu}_{ab}({p^2},m)
   =
   \delta_{ab} \x \left(p^\mu p^\nu - g^{\mu\nu} p^2\right) \; \Pi({p^2},m)
$$
\vspace{-3mm}
$$
   \Pi(\green{p^2},m)
   =
   \int \D^n l \; F(p,l,m)
$$
  \item {\bf Arguments}: {\em from vectors to scalars}
        $$ F(p,l,m) \;\to\; F(\blue{l^2},\; \blue{l \!\cdot\! p}, \; \green{p^2}, \; m) $$
  \item In general, the number of \blue{scalar integration variables} is given by
        $$N(L,E) = \frac{L\,(L+1)}{2} + L\,E \quad \sim \quad \mathcal{O}(L^2) \;\;\leftarrow\;\; \parbox{50mm}{another source of growing complexity at higher orders}$$%
        where $E$ -- number of {\em external momenta}, $L$ -- number of {\em loop momenta}
  \bi
    \item 1-loop propagator: $\blue{N(1,1)=2}$
    \item 4-loop propagator: $N(4,1)=14$ (ask Jos Vermaseren about details)
  \ei
\ei
\newpage


\titleboxx{I.1}{\Large Integration-by-Parts Reduction}
\bi
  \item The problem contains two denominators 
        $$\quad D_1 = l^2-m^2 \quad \quad D_2 = (l-p)^2-m^2  $$%
        which map into our integration invariants in a unique way 
        $$ F(p,l,m) \;\to\; F(\blue{l^2},\; \blue{l \!\cdot\! p}, \; p^2, \; m) \to F(\blue{D_1}, \blue{D_2}, p^2,m) $$
  \item One integral family
        $$F(n_1, n_2) = \int \D^n l \frac{1}{D_1^{n_1} D_2^{n_2}}$$
\ei

\vspace{4mm}
\begin{Verbatim}[fontsize=\small,frame=lines, label=\fbox{Mathematica}]

  <<LiteRed`

  SetDim[n];
  Declare[{m2}, Number, {l,p}, Vector];
  NewBasis[$b, {sp[l]-m2, sp[l-p]-m2}, {l}, Directory->"b.ibp"];

  GenerateIBP[$b];
  AnalyzeSectors[$b];
  FindSymmetries[$b];

\end{Verbatim}
\newpage


\titleboxx{I.1}{\Large Master Integrals}
\begin{Verbatim}[fontsize=\small,frame=lines, label=\fbox{Mathematica}]

  SolvejSectors /@ UniqueSectors[$b]

  MIs[$b]

  > {j[$b,0,1], j[$b,1,1]}
\end{Verbatim}
\bi
  \item We obtain two master integrals
        $$F_1 = F(0,1) =  \int \D^n l \frac{1}{(l-p)^2-m^2} \qquad F_2 = F(1,1) =  \int \D^n l \frac{1}{\big(l^2-m^2\big)\x\big((l-p)^2-m^2\big)}$$
  \item Any other integral is a linear combination of only these two, e.g.,
        $$F(2,1) = \frac{n-2}{2 m^2 (p^2-4m^2)} F_1 + \frac{n-3}{p^2-4m^2} F_2$$
  \item We can check that since we can do $l\to l+p$ transformation $$F(0,1) = F(1,0)$$
\ei
\newpage


\titleboxx{I.2}{\Large Differential Equations}

\vs
\begin{Verbatim}[fontsize=\small,frame=lines, label=\fbox{Mathematica}]

$ds = Dinv[#,sp[p,p]]& /@ MIs[$b] // IBPReduce;
$ode = Coefficient[#, MIs[$b]]& /@ $ds;
\end{Verbatim}
\vs

\bi
  \item This code produces a system of differential equations
  $$
    \begin{aligned}
     \frac{\D F_1}{\D p^2} & = 0
     \\
     \frac{\D F_2}{\D p^2} & = \frac{2-2\eps}{p^2\x(p^2-4m^2)} F_1 + \frac{2m^2-\eps p^2}{p^2\x(p^2-4m^2)} F_2
    \end{aligned}
  $$
  where we work in $n=4-2\eps$ space-time dimensions
\ei
\vs

This system is simple and we could solve it right away using {\em <your favourite>} method.
Today, I want to demonstrate you how this and many other systems can be solved throug using their {\bf $\boldsymbol{\eps}$-form}.
As you will see this is a highly automated task.

\vs
\fbox{Exercise 1a.} \\
Derive another system of differential equations, but this time in~$m^2$. (Hint: instead of \code{Dinv} use \code{Fromj}, \code{D}, and \code{Toj} functions ).
\newpage


\titlebox{Epsilon Form}
\bi
  \item Classical Notation
$$
\begin{aligned}
 \frac{\D F_1}{\D x} & = A_{11}(x,\eps) F_1 + A_{12}(x,\eps) F_2
 \\
 \frac{\D F_2}{\D x} & = A_{21}(x,\eps) F_1 + A_{22}(x,\eps) F_2
\end{aligned}
$$
  \item Matrix Notation
$$
\frac{\D \bar{F}}{\D x} = A(x,\eps) \x \bar{F}
\qquad \text{where}\qquad
A=
\begin{pmatrix}
  A_{11}(x,\eps) & A_{12}(x,\eps) \\
  A_{21}(x,\eps) & A_{22}(x,\eps)
 \end{pmatrix}
\quad \text{and}\quad
\bar{F}=
\begin{pmatrix}
  F_1 \\
  F_2
 \end{pmatrix}
$$
\ei
It is very convenient (as observed by J.Henn) to have our system in the epsilon form
$$
\frac{\D G}{\D x} = \eps \x B(x) \x G
$$%
since in this case we can easily find the solution to any order in $\eps$ parameter, as we will see on the next slide.
\newpage


\titlebox{A few words on \fuchsia}
\fuchsia \ldots
\newpage


\titleboxx{I.3}{\Large Epsilon Form by \fuchsia}
Let us introduce a new variable $y$, such that
$$p^2 = -4m^2\frac{y^2}{1-y^2}$$%
The new equations look as
  $$
    \begin{aligned}
     \frac{\D F_1}{\D y} & = 0
     \\
     \frac{\D F_2}{\D y} & = \frac{1-\eps}{y \x m^2} F_1 + \left( \frac{\eps}{1-y} - \frac{\eps}{1+y} - \frac{1}{y} \right) F_2
    \end{aligned}
  $$

With the help of \fuchsia we find a new basis $G_1$, $G_2$ given by the system
  $$
    \begin{aligned}
     F_1 & = \frac{4(1-2\eps)}{3(1-\eps)} G_1
     \\
     F_2 & = \frac{4}{3m^2} G_1 - \frac{2}{y} G_2
    \end{aligned}
  $$%
For this basis the differential equations are the epsilon form
  $$
    \begin{aligned}
     \frac{\D G_1}{\D y} & = 0
     \\
     \frac{\D G_2}{\D y} & = \frac{2}{3m^2}\left( \frac{\eps}{1+y}+\frac{\eps}{1-y} \right) G_1 - \left(\frac{\eps}{1-y}-\frac{\eps}{1+y}\right) G_2
    \end{aligned}
  $$
\newpage


\titlebox{Solutions}
We are looking for the solution of a given system of ordinary differential equations in the epsilon form
$$
\frac{\D G}{\D x} = \eps \x B(x) \x G
$$%
as a Laurent series in $\eps$
$$
  \bbox{
  G(x,\eps) = G_0(x) + G_1(x) \x \eps + G_2(x) \x \eps^2 + \ldots}
$$%
Let us put this "solution" into the initial equation
$$
\frac{\D G_0}{\D x} + \frac{\D G_1}{\D x} \eps + \frac{\D G_2}{\D x} \eps^2 + \ldots = \eps \x B(x) \x G_0 + \eps^2 \x B(x) \x G_1
$$%
we get
$$
 \frac{\D G_0}{\D x} = 0,
 \qquad
 \frac{\D G_1}{\D x} = B(x) G_0,
 \qquad
 \frac{\D G_2}{\D x} = B(x) G_1
 \qquad
 \ldots
 \qquad
 \frac{\D G_{n}}{\D x} = B(x) G_{n-1}
$$%
This system can be easily solved (as promised)
$$
 G_0 = C_0,
 \qquad
 G_1 = C_1 + \int \D x B(x) C_0,
 \qquad
 G_2 = C_2 + \int \D x B(x) \left( C_1 + \int \D x B(x) C_0\right)
 \qquad
 \ldots
$$
$$
\bbox{
 G_n(x) = C_n + \int \D x B(x) G_{n-1}
}
$$
\newpage


\titlebox{Solutions}

My implementation of the solution algorithm, which I use to get results for the next slide.
\vs

\begin{Verbatim}[fontsize=\small,frame=lines, label=\fbox{Mathematica}]

  SolveODE[m_, x_, ep_, n_, c_] := Module[
    {$i, $j, $n, $sol, $sol0, $sol1},

    $n = Length[m];
    $sol[0] = Table[c[$j,0], {$j,1,$n}];
    For[$i=1, $i<=n, $i++,
      $sol0 = Table[c[$j,$i], {$j,1,$n}];
      $sol1 = Integrate[Dot[#,$sol[$i-1]],x]& /@ m;
      $sol[$i] = $sol0 + $sol1;
    ];
    Sum[ep^$i*$sol[$i], {$i,0,n}]
  ];

\end{Verbatim}
\newpage


\titleboxx{I.4}{\Large Solutions}
\bi
  \item Master \#1
    $$ F_1(y,m^2) = \frac{4}{3}C_1^{(0)} +  \frac{4}{3}\left(C_1^{(1)}-C_1^{(0)}\right)\eps + \ldots $$
  \item Master \#2
    $$
      F_2(y,m^2)
      =
      \frac{4C_1^{(0)}}{3m^2}
      -
      \frac{C_2^{(0)}}{y}
      +
      \frac{\eps}{3m^2y}
      \left(
          4yC_1^{(1)}
        - 6m^2C_2^{(1)}
        + \big(4 C_1^{(1)}
        - 6 m^2 C_2^{(0)}\big) \ln\left(\frac{1-y}{1+y}\right)
      \right)
    $$
  \item Finally, we need to find unknown integration constants whcih are \\ functions of $m^2$ and~$\eps$, i.e.
        $$ C_1^{(0)}(m^2,\eps), \quad  C_1^{(1)}(m^2,\eps), \quad \ldots $$
        $$ C_2^{(0)}(m^2,\eps), \quad  C_2^{(1)}(m^2,\eps), \quad \ldots $$
\ei
\newpage


\titleboxx{I.5}{\Large Integration Constants \#1}
    Master \#1 (from \fuchsia)
    $$ F_1(y,m^2) = \frac{4}{3}C_1^{(0)} +  \frac{4}{3}\left(C_1^{(1)}-C_1^{(0)}\right)\eps + \ldots $$

    Closed-form solution from the literature (see Smirnov's book)
    $$ F(0,n) = (-1)^n \frac{\Gamma(n-2+\eps)}{\Gamma(n)} \x (m^2)^{2-\eps-n} $$
    $$
      F_1(y,m^2)
       = F(0,1) 
       = \frac{m^2}{\eps} + m^2 \left(1- \gamma_E - \ln{m^2} \right) + \ldots
    $$

Result \#1
    $$ \bbox{C_1^{(0)} = \frac{3m^2}{4\eps} \quad C_1^{(1)} = \frac{3m^2 \left(2- \gamma_E - \ln{m^2} \right)}{4\eps}}$$
\newpage


\titleboxx{I.5}{\Large Integration Constants \#2}
  Result \#1
    $$ C_1^{(0)} = \frac{m^2}{\eps} \quad C_1^{(1)} = \frac{m^2 \left(1- \gamma_E - \ln{m^2} \right)}{\eps}$$

  Master \#2 (with Result \#1 substituted)
    $$
      F_2(y,m^2)
      =
      \frac{1}{\eps}
      +
      \frac{2y - \gamma_E\x y - 2 C_2^{(0)} - y\x \ln{m^2} + \ln\left(\frac{1-y}{1+y}\right)}{y}
      +
      \ldots
    $$

  We require that at the limit $y\to 0$ ($p^2\to 0$) our result is regular. This leads to the solution
$$ C_2^{0} = 0 $$

Result \#2
$$\bbox{
  F_2(y,m^2)
  =
  \frac{1}{\eps}
  +
 2 - \gamma_E\x - \ln{m^2} + \frac{1}{y}\ln\left(\frac{1-y}{1+y}\right)
  +
  \ldots}
$$

This is in agreement with T.Riemann Monday's lecture!
\newpage


\titleboxx{Example I}{\Large \bf Conclusions}
\newpage


\titlebox{II.1 Splitting Functions from $e^+ e^-$-annihilation}
In this example, I will show how to calculate a gluon-quark splitting function
$$P_{gq} = \frac{1 + (1-x)^2}{x}$$%
Using this technique you will be able to calculate remaining splitting functions $P_{qq}$, $P_{qg}$, and $P_{gg}$ as well as higher-order corrections to these quantities.

\begin{center}
\begin{axopicture}(320,80)
  \SetScale{0.75}
  \Line[arrow](60,50)(20,90)
  \Line[arrow](20,10)(60,50)
  \Vertex(60,50){2}
  \Photon(60,50)(120,50){5}{4}
  \Vertex(120,50){2}
  \Line[arrow](120,50)(160,90)
  \Line[arrow](160,10)(120,50)
  \Gluon(130,60)(160,60){-4}{3}
  \Vertex(130,60){2}
  \Text(90,65){$q$}
  \Text(10,90){$q_1$}
  \Text(10,10){$q_2$}
  \Text(170,90){$p_1$}
  \Text(170,60){$p_3$}
  \Text(170,10){$p_2$}

  \SetOffset(180,0)
  \Line[arrow](60,50)(20,90)
  \Line[arrow](20,10)(60,50)
  \Vertex(60,50){2}
  \Photon(60,50)(120,50){5}{4}
  \Vertex(120,50){2}
  \Line[arrow](120,50)(160,90)
  \Line[arrow](160,10)(120,50)
  \Gluon(130,40)(160,40){-4}{3}
  \Vertex(130,40){2}
  \Text(90,65){$q$}
  \Text(10,90){$q_1$}
  \Text(10,10){$q_2$}
  \Text(170,90){$p_1$}
  \Text(170,40){$p_3$}
  \Text(170,10){$p_2$}
\end{axopicture}
\end{center}
$$e^+(q_1) + e^-(q_2) \to q(p_1) + \bar{q}(p_2) + g(p_3)$$

We employ a mass-factorization theorem 
$$
 \frac{\D\sigma_i}{\D x} = \frac{P_{iq}}{\eps} + a_i + b_i \eps + \ldots
$$%
where $q=q_1+q_2$ and
$$
\frac{\D\sigma_i}{\D x} = \int \D^n p_1 \delta(p_1^2) \x \D^n p_2 \delta(p_2^2) \x \D^n p_3 \delta(p_3^2) \x \delta\left(x-\frac{2 q\!\cdot\!p_i}{q^2}\right) \x \sigma(q_1,q_2,p_1,p_2,p_3)
$$
\newpage


\titlebox{II.2 Phase-Space Integrals}
By their structure phase-space integrals are very similar to loop integrals (compare to the one-loop propagator from Example I), except that we apply on-shell conditions $\delta(p_i^2)$ to the cut lines as shown in the following {\bf cut graph}

\begin{center}
\begin{axopicture}(220,80)
  \SetScale{0.75}
  \Photon(80,50)(120,50){5}{3}
  \Vertex(120,50){2}
  \Line[arrow](160,90)(120,50)
  \Line[arrow](120,50)(160,10)
  \Gluon(130,60)(190,60){-4}{6}
  \Vertex(130,60){2}
  \Vertex(190,60){2}
  \Line[arrow](200,50)(160,90)
  \Line[arrow](160,10)(200,50)
  \Vertex(200,50){2}
  \Photon(200,50)(240,50){5}{3}
  \DashLine(160,5)(160,95){6}
\end{axopicture}
\end{center}

$$
\frac{\D\sigma_g}{\D x} = \int \D^n p_1 \delta(p_1^2) \x \D^n p_2 \delta(p_2^2) \x \D^n p_3 \delta(p_3^2) \x \delta\left(x-\frac{2 q\!\cdot\!p_3}{q^2}\right) \x \sigma(q_1,q_2,p_1,p_2,p_3)
$$
where
$$
\sigma(q_1,q_2,p_1,p_2,p_3)
=
N \frac{(\spr{p_1}{q_1})^2+(\spr{p_2}{q_1})^2+(\spr{p_1}{q_2})^2+(\spr{p_2}{q_2})^2}{\spr{p_1}{p_3}\x\spr{p_2}{p_3}}
$$
This integration is equivalent to the 2-loop propagator, since we can eliminate one of the integration momenta using momentum conservation
$$ q_1 + q_2 = p_1 + p_2 + p_3 $$
\newpage


\titlebox{II.3 Integration-by-Parts}

In order to integrate the cross-section we need a new IBP basis. Let us define one as
\vs

\begin{Verbatim}[fontsize=\small,frame=lines, label=\fbox{Mathematica}]

NewBasis[$a, {sp[p1], sp[p3], sp[q1+q2-p1-p3], s*x-2sp[q1+q2,p3], sp[p1,p3]},
             {p1, p3}, Append -> True];

GenerateIBP[$a];
AnalyzeSectors[$a, {___,0,0}, CutDs -> {1,1,1,1,0,0,0}];
FindSymmetries[];

SolvejRules /@ UniqueSectors[$a];

\end{Verbatim}
\vs
Note additional arguments in \code{AnalyzeSectors} routine:
\bi
  \item in \code{\{\_\_\_,0,0\}} 0's represent invariants which appear in numerators only
  \item in \code{CutDs -> \{1,1,1,1,0,0,0\}} 1's represent "cut" propagators. It means that all integrals with at least one non-positive indices in these places vanish.
\ei

\vs
We get only one master integral
$$
F_1(x,\eps) = \int \D^n p_1 \delta(p_1^2) \x \D^n p_2 \delta(p_2^2) \x \D^n p_3 \delta(p_3^2) \x \delta\left(x-\frac{2 q\!\cdot\!p_3}{q^2}\right)
$$

\newpage


\titlebox{Partial Fractioning}
Given a set of denominators, being a {\em linear combination} of the kinematic invariants $s_{ij}$, make a partial fraction such that
$$
  \frac{1}{D_1\ldots D_n} \to \frac{a_1}{D_2 \ldots D_n} +  \frac{a_2}{D_1 D_3 \ldots D_n} + \ldots + \frac{a_n}{D_1 \ldots D_{n-1}}
$$

All we need is to solve a linear system of equations
$$ a_1 D_1 + \ldots + a_n D_n = N $$
where the coefficient in front of every $s_{ij}$ is zero and $N$ is some number.

\vs
In particaulr, for
$$
  A = \frac{1}{(x+1)(y+1)(x+y+1)}
$$
we write down
$$
  (a_1 + a_3) x + (a_2 + a_3) y + a_1 + a_2 + a_3 = N
$$
the solution is
$$
  a_1 = - a_3 \quad  a_2 = - a_3 \quad N = - a_3
$$
which gives
$$
  A = \frac{1}{(y+1)(x+y+1)} + \frac{1}{(x+1)(x+y+1)} - \frac{1}{(x+1)(y+1)}
$$
\newpage


\titlebox{II.3 Integration-by-Parts}

Now we can convert the initial cross-section into the $j$-form and make IBP reduction.
\begin{Verbatim}[fontsize=\small,frame=lines, label=\fbox{Mathematica}]

$PS2 = x / (sp[p1] sp[p3] sp[q1+q2-p1-p3] (s*x-2*sp[q1+q2,p3]));

$jM2 = Toj[$a, $PS2 $M2];

$jM2 = $jM2 // IBPReduce

$Pgq = Series[$jM2 /. {m -> 4 - 2 eps}, {eps, 0, -1}]

\end{Verbatim}
\vs
This gives us
$$ P_{gq} \sim \frac{2-2x+x^2}{\color{red}x^2} \x F_1(x)  $$%
which contains one $x$ factor more in the denominator.

\vs\vs
Maybe $F_1(x) \sim x$? Let us check this\ldots

\newpage


\titlebox{II.4 Differential Equations}

\vs
\begin{Verbatim}[fontsize=\small,frame=lines, label=\fbox{Mathematica}]

$F1 = j[$a, 1, 1, 1, 1, 0, 0, 0];

$dF1 = Toj[$a, D[Fromj[$F1], x]] // IBPReduce;

\end{Verbatim}
\vs
This code produces the following equation
$$
  \frac{\D F_1}{\D x} = \left( \frac{\eps}{1-x}+\frac{1-2\eps}{x} \right) F_1
$$

Of course we could use \fuchsia and find the $\eps$-form, but we can solve this in a closed form
$$
  F_1 = C(\eps) \x (1-x)^\eps \x x^{1-2\eps}
$$
which confirms our assumption from the previous slide.

The final result is
$$ \bbox{P_{gq} \sim \frac{2-2x+x^2}{x}}$$%

\newpage



\titlebox{Holonomic Functions}
A function $f=f(x)$ is called {\em holonomic} if there exist polynomials $a_n(x),\x\ldots,\x a_0(x)$ such that
$$
  \label{eq:ode}
  a_n(x) f^{(n)} - a_{n-1}(x)\x f^{(n-1)} - \ldots - a_0(x) \x f = 0
$$%
holds for all $x$. Hence, the holonomic function is uniquely defined by
\bi
  \item the differential equation
  \item a number of initial values $f(x_0),\x f'(x_0),\x \ldots,\x f^{(n-1)}(x_0)$
\ei

\vs
\underline{Examples} of holonomic functions:
\bi
  \item {\em all algebraic functions}
  \item \fbox{\em Generalized Hypergeometric functions}
    \bi
      \item {\em polylogarythms}
      \item {\em Elliptic functions}
    \ei
  \item {\em Bessel functions}
  \item {\em Airy functions}
  \item {\em Legendre and Chebyshev polynomials}
  \item {\em Heun functions}
  \item and many others that have no name and no closed form
\ei

\vs\vs
For more details, see the lecture by S.Moch at CAPP17 and  S.Kauers at CAPP11.
\newpage

\titlebox{Holonomic Functions}
A function $f=f(x)$ is called {\em holonomic} if there exist polynomials $a_n(x),\x\ldots,\x a_0(x)$ such that
$$
  a_n(x) f^{(n)} - a_{n-1}(x)\x f^{(n-1)} - \ldots - a_0(x) \x f = 0
$$%
holds for all $x$. Hence, the holonomic function is uniquely defined by
\bi
  \item the differential equation
  \item a number of initial values $f(x_0),\x f'(x_0),\x \ldots,\x f^{(n-1)}(x_0)$
\ei

\vspace{1cm}
\fbox{Conclusion}
\vs
\bn
\item With the help of polynomials and differential equations we can define a large set of complicated functions which have no closed form and usually have a non-trivial integration representation.
\item These functions appear in the Feynman integrals.
\item In some situation it is much easier to solve ODE than to integrate directly.
\en
\newpage


\titlebox{$n^{\text{th}}$-order ODE vs $n\times n$ system}

We can easily rewrite a $n^{\text{th}}$-order linear ODE given by
\be
  \label{eq:ode}
  y^{(n)} - a_1(x)\x y^{(n-1)} - \ldots - a_n(x) \x y = 0
\ee
as an $n\times n$ system of the form 
\be \frac{\D \bar{y}}{\D x} = A(x) \x \bar{y} \ee
where
\be
A(x) =
\begin{bmatrix}
  0 & 1 & \cdots & 0 & 0
  \\
  0 & 0 & \cdots & 0 & 0
  \\
  \vdots & \vdots & \ddots & \vdots & \vdots
  \\
  0 & 0 & \cdots & 0 & 1
  \\
  a_n(x) & a_{n-1}(x) & \cdots & a_2(x) & a_1(x)
\end{bmatrix}
\quad \text{and} \quad
\bar{y} =
\begin{pmatrix}
  y \\ y' \\ \vdots \\ y^{(n-2)} \\ y^{(n-1)}
 \end{pmatrix}
\ee
However, the inverse opperation is not as easy anymore.
\newpage


\titleboxx{Example II}{Expansion of Hypergeometric Functions}
The Generalized Hypergeometric Function
\be
 _{p+1}F_p
\left(
\begin{matrix}
  a_1,\x a_2,\x \ldots,\x a_{p+1}
  \\
  b_1,\x b_2,\x \ldots,\x b_{p}
\end{matrix}
;\x x
\right)
=
\prod_{i=1}^{p}\frac{\Gamma(b_i)}{\Gamma(a_i)\Gamma(b_i-a_i)}
\int_0^1 \frac{t_i^{a_i-1} (1-t_i)^{b_i-a_i-1}}{(1-x\x t_1\ldots t_{p})^{a_{p+1}}} \D t_i
\ee
is a solution to the differential equation
\be
\left[
 D \x (D+b_1-1)\x\cdots\x(D+b_p-1)\x - \x x\x(D+a_1)\x\cdots\x(D+a_{p+1})
\right] \x y = 0
\ee
where
$$D = x\frac{\D}{\D x}$$

\vs
\fbox{Exercise 2.}
Using your favourite CAS write a routine which for a given Generalized Hypergeometric Function, defined by the list $\{a_1,\ldots,a_{p+1},b_1,\ldots,b_p\}$, it returns a corresponding ODE, defined by the list $\{a_1(x),\ldots,a_p(x)\}$, in accordance with notation of eq.~\eqref{eq:ode}.

\newpage


\titlebox{A Word of Advice}
\bi
  \item Before computing, first think of how to check your results
  \item Read arXiv.org daily
  \bi
    \item subscribe to the mailing list
    \item e-mail To: physics@arxiv.org, Subject: subscribe
  \ei
  \item Search with scholar.google.com
  \item Do not be afraid to write e-mails
  \item Discuss your ideas
\ei
\newpage

\titlebox{Reading List}
\bi \item {\em "Feynman Integral Calculus"} by V.~Smirnov
    \item {\em "Lectures on Differential Equations for Feynman Integrals"} by J.~Henn
    \item {\em "Formal Power Series and Linear Systems of Meromorphic Ordinary Differential Equations"} by W.~Balser
    \item {\em "Computer Algebra in Particle Physics"} by S.~Weinzierl
    \item {\em "Introduction to Loop Calculations"} by G.~Heinrich
    \item {\em "Structure and Interpretation of Computer Programs"}\\by H.~Abelson and G.~Sussman with J.~Sussman
    \ei
\newpage

\titlebox{Summary}

\end{document}
